\section*{Capítulo \ref{ch:g1:shapes} --- Formas e espaço}

\begin{solution}{exo:g1:shapes:1}
$3$ cantos: um triângulo. Redonda, sem cantos: um círculo. $4$ lados
iguais: um quadrado.
\end{solution}

\begin{solution}{exo:g1:shapes:2}
As respostas variam: uma porta e um livro são retângulos; um prato
ou um relógio é um círculo; uma placa de advertência é um triângulo;
um azulejo pode ser um quadrado.
\end{solution}

\begin{solution}{exo:g1:shapes:3}
Falso: um quadrado apoiado num canto continua com quatro lados
iguais e quatro cantos --- girar não muda a forma.
\end{solution}

\begin{solution}{exo:g1:shapes:4}
Desenhos na malha: o quadrado tem $3 \times 3$ quadradinhos, o
retângulo $5 \times 2$, e o triângulo usa uma linha inclinada da
malha como um dos lados.
\end{solution}

\begin{solution}{exo:g1:shapes:5}
As respostas variam. Por exemplo: a bola está à esquerda, o carrinho
está entre a bola e a boneca, e a boneca está à direita.
\end{solution}

\begin{solution}{exo:g1:shapes:6}
A cruz é desenhada dentro do círculo, a estrela fora, e o ponto em
cima da própria linha do círculo.
\end{solution}

\begin{solution}{exo:g1:shapes:7}
Seguir $\to\ \uparrow\ \to\ \uparrow\ \to$ desloca $3$ quadradinhos
para a direita e $2$ para cima. O caminho de volta é
$\downarrow\ \downarrow$ e depois
$\leftarrow\ \leftarrow\ \leftarrow$ (ou qualquer caminho que volte
à saída).
\end{solution}

\begin{solution}{exo:g1:shapes:8}
Dois triângulos e um quadrado: $3 + 3 + 4 = 10$ cantos. Um círculo e
dois retângulos: $0 + 4 + 4 = 8$ cantos.
\end{solution}

\begin{solution}{exo:g1:shapes:9}
Com $12$ palitos dá para fazer um quadrado com $3$ palitos em cada
lado ($3 + 3 + 3 + 3 = 12$). Com $10$ palitos não dá: $10$ não se
reparte em $4$ partes inteiras iguais (cada lado ficaria com $2$
palitos e meio).
\end{solution}
